% from main.tex
%%%%%%%%%%  scan idx 94  |  folio 81  %%%%%%%%%%
Pseudo-cohérence et perfection sont les deux notions de finitude fondamentales avec lesquelles nous travaillerons dans ce Séminaire. Les §§~1 à 5 de l'Exposé I sont consacrés à leur définition et à l'établissement de leurs propriétés élémentaires de stabilité. Pour ceci, on utilise seulement deux ou trois propriétés locales de la catégorie des Modules libres de type fini vis-à-vis de la catégorie de tous les Modules. Aussi s'avère-t-il commode~-- et utile, notamment dans l'Exp.\ II~-- d'axiomatiser la situation, en introduisant une notion de pseudo-cohérence (resp.\ de perfection) dans une catégorie fibrée sur un site relativement à une sous-catégorie fibrée vérifiant des conditions convenables.

Les §§~6 à 8 de l'Exp.\ I généralisent aux complexes parfaits un certain nombre de notions bien connues pour les faisceaux localement libres de type fini : rang, dualité, trace d'un endomorphisme.

Dans l'Exp.\ II, on examine le problème de «~l'existence de résolutions globales~» : à quelles conditions, sur un topos annelé \((X,\mathcal O_X)\), un complexe pseudo-cohérent (resp.\ parfait) est-il \emph{globalement} isomorphe dans \(\Db(X)\) à un complexe de Modules localement libres de type fini? L'importance de cette question réside essentiellement dans le fait qu'on ne sait pas, pour le moment, généraliser aux complexes parfaits certaines constructions usuelles sur les Modules localement libres : puissances extérieures et symétriques, classes de Chern (pour plus de détails à ce sujet, voir Exp.\ XIV). C'est pourquoi il est commode d'avoir à sa disposition des critères suffisants maniables pour l'existence de résolutions globales, ce qui permet de ramener certaines questions sur les complexes pseudo-cohérents (resp.\ parfaits) à des questions analogues sur les faisceaux localement libres de type fini. De tels critères
